\section{Розділ~\ref{sec:serocomputer}. \nt{SERO}-нейрони}

Властивості самих \nt{SERO}-нейронів (ритм, автогальмування, знак за рецептором, підсистеми) виміряно прямо; обчислення розділу --- регулятор, фільтр, логіка --- випливають із його рівнянь; $\tau_c$, коефіцієнт дифузії серотоніну і кількість місць викиду --- оцінки; роль петлі як регулятора рівня в мозку --- гіпотеза (у ядрі шва автогальмування точкове і дає лише короткі паузи), а роль \nt{SERO} як горизонту --- гіпотеза, на користь якої свідчать поведінкові досліди.

{\footnotesize
\setlength{\tabcolsep}{3pt}\renewcommand{\arraystretch}{1.15}
\begin{longtable}{>{\raggedright}p{0.5cm}>{\raggedright}p{4.0cm}>{\raggedright}p{5.6cm}>{\raggedright}p{2.3cm}>{\raggedright\arraybackslash}p{1.7cm}}
\toprule
№ & Твердження & Дослід і що виміряно & Джерело & Статус \\
\midrule\endfirsthead
\toprule
№ & Твердження & Дослід і що виміряно & Джерело & Статус \\
\midrule\endhead
1 & Знак дії серотоніну вибирає рецептор отримувача (твердження~\ref{prop:receptor}) & Рецептори серотоніну поділено на родини за фармакологією і механізмом: 5-HT$_{1A}$ через G-білок відкриває калієві канали і гальмує клітину, 5-HT$_{2A}$ та іонотропний 5-HT$_3$ збуджують. Один медіатор дає протилежні відповіді в різних клітинах. & \cite{BarnesSharp1999} & виміряно \\
2 & \nt{SERO}-нейрон у стані неспання рівно видає кілька імпульсів на секунду & Запис нейронів дорсального ядра шва у вільно рухливих котів: повільний ритмічний розряд $2.8\pm0.2$ імп/с у спокійному неспанні; у повільному сні частота падає на $34$--$68\,\%$, у швидкому --- на $98\,\%$. Під наркозом у щурів --- $1.7\pm0.5$ імп/с, дуже регулярно. & \cite{TrulsonJacobs1979, GagnonParent2014, JacobsAzmitia1992} & виміряно \\
3 & \nt{SERO}-нейрон --- ритмоводій: поштовх на кожен імпульс не потрібен, але потрібен рівний фоновий приток (у зрізі --- норадреналін через $\alpha_1$-рецептори, в організмі --- від \nt{NORA}) & У зрізі мозку клітини розряджаються повільно і рівно, після кожного імпульсу --- слідова гіперполяризація на $200$--$800\ms$. Спонтанно активних клітин у зрізі менше, ніж в організмі: рівному ритму потрібен тонічний вхід норадреналіну через $\alpha_1$-рецептори, блокада яких його гасить. & \cite{VandermaelenAghajanian1983} & виміряно \\
4 & Майже всі рецептори метаботропні; відповідь повільна: наростає і спадає в межах приблизно секунди & Усі родини рецепторів, крім 5-HT$_3$, пов'язані з G-білками. У зрізі викликаний викид серотоніну дає через 5-HT$_{1A}$ гальмівний струм, який наростає і спадає в межах секунди. & \cite{BarnesSharp1999, CourtneyFord2016} & виміряно \\
5 & В областях призначення серотонін виходить в об'єм, а не лише в щілину; у самому ядрі шва гальмування сусідів іде точка-точка & Швидка циклічна вольтамперометрія у зрізах: викид на один імпульс однаковий в області із синапсами і в ядрі шва, де синапсів майже немає; не залежить від блокади захоплення і рецепторів; молекул на закінчення набагато менше, ніж переносників і рецепторів, а позаклітинна концентрація збігається зі спорідненістю головного рецептора. У ядрі шва гальмування через 5-HT$_{1A}$ іде точка-точка: переносник не дає серотоніну накопичитися поза щілиною. & \cite{Bunin1998, CourtneyFord2016} & виміряно \\
6 & Концентрація за~\eqref{eq:seroneuron} спадає через відкачування; $\tau_c$ порядку секунди --- оцінка & Вольтамперометрія в живому мозку: позаклітинний серотонін обмежують насамперед зворотне захоплення і руйнування, а не синтез. Прямого вимірювання $\tau_c$ у цитованих роботах немає; порядок величини --- оцінка розділу. & \cite{Hashemi2012serotonin} & виміряно; $\tau_c$ --- оцінка \\
7 & НЕ і АБО-НЕ на одному ритмоводії; повнота \nt{SERO}-логіки (твердження~\ref{prop:serocomplete}, рис.~\ref{fig:serogates}) & Випливає з~\eqref{eq:seroneuron}: гальмований ритмоводій --- АБО-НЕ, а АБО-НЕ функціонально повний. Досліду, де з \nt{SERO}-нейронів зібрано логіку, немає; умова окремих об'ємів у мозку не виконується. & виведено в розділі & теорія \\
8 & Серотонін гальмує сам \nt{SERO}-нейрон через рецептори на ньому & У щурів під наркозом агоністи 5-HT$_{1A}$ внутрішньовенно та іонофоретично гальмують розряд нейронів дорсального шва з тією самою залежністю доза--відповідь, що й сам серотонін; агоністи 5-HT$_{1B}$ майже не діють. У зрізі ендогенний серотонін сусідніх клітин дає через 5-HT$_{1A}$ струм і коротку паузу в розряді, не змінюючи середньої частоти. & \cite{Sprouse1987, CourtneyFord2016} & виміряно \\
9 & У моделі петля через спільний об'єм --- регулятор рівня: розкид і стала часу зменшуються в $1+G$ разів~\eqref{eq:feedback}; що петля так працює в мозку, --- гіпотеза & Класична формула підсилювача з від'ємним зворотним зв'язком; у розділі виведена заново. Підсилення петлі $G$ у \nt{SERO}-системи не виміряно. Виміряно: у зрізі автогальмування іде точка-точка, дає коротку паузу і не змінює середньої частоти. & \cite{Black1934feedback, CourtneyFord2016} & теорія; у мозку --- гіпотеза \\
10 & Концентрація --- ковзне середнє частоти, фільтр нижніх частот~\eqref{eq:lowpass}, рис.~\ref{fig:lowpass} & Розв'язок лінійного рівняння~\eqref{eq:seroneuron}; рисунок --- розрахунок за цією формулою при $\tau_c=1\,\text{с}$. Окремої моделі в \texttt{models/} немає. & виведено в розділі & теорія \\
11 & Звивистість щілин сповільнює дифузію малої молекули приблизно в $2.6$ раза & Метод точкового джерела: іонофорез тетраметиламонію і запис його концентрації іоноселективним електродом у зрізах і в живому мозку. Звивистість $\lambda\approx1.6$, тобто ефективний $D$ менший за вільний у $\lambda^2\approx2.6$ раза; частка міжклітинного об'єму близько $0.2$. Коефіцієнт дифузії самого серотоніну в тканині в цих роботах не вимірювали; у розділі він --- оцінка. & \cite{Sykova2008} & виміряно \\
12 & Довжина незалежного «проводу» $\ell\sim\sqrt{D\tau}\approx20$~мкм~\eqref{eq:diffusionlength}; кількість проводів обмежена кількістю таких областей (твердження~\ref{prop:volumewires}) & Закон дифузії і підстановка оцінок $D$ і $\tau$ (рядки~6 і~11); твердження~\ref{prop:volumewires} --- наслідок. Прямого досліду, який рахує незалежні канали об'ємної передачі, немає. & виведено в розділі & теорія \\
13 & Вихід одного \nt{SERO}-нейрона розкинутий по величезних ділянках мозку; місць викиду, за оцінкою, $\sim10^5$ & Повна реконструкція окремих нейронів дорсального шва в щура: усі висхідні аксони сильно розгалужуються, загальна довжина аксона однієї клітини до $18.7$~см, гілки однієї клітини сягають стріатума, префронтальної кори і мигдалеподібного тіла. Кількість місць викиду на клітину прямо не підраховано. & \cite{GagnonParent2014} & виміряно; $10^5$ --- оцінка \\
14 & \nt{SERO}-нейрони поділяються на кілька підсистем із різними виходами і відповідями & Вірусно-генетичне мічення в мишей: нейрони, що йдуть до мигдалеподібного тіла і до лобової кори, майже не перетинаються за розгалуженням, отримують різні входи і протилежно відповідають на неприємний стимул; увімкнення і вимкнення кожної групи по-різному змінює поведінку. & \cite{Ren2018} & виміряно \\
15 & Умова терпіння $D<\tau_\gamma\ln(R/r_0)$~\eqref{eq:patience} за експоненційного знецінення; за виміряного, гіперболічного, --- $D<\tau_\gamma(R/r_0-1)$ & Випливає зі знецінення $e^{-D/\tau_\gamma}$ або $1/(1+D/\tau_\gamma)$; виведено в розділі. Мавпи, вибір на затримках у секунди: знецінення ближче до гіперболи, ніж до експоненти; відповідь \nt{DOPA}-нейронів на провісник відкладеної винагороди спадає із затримкою теж за гіперболою. & виведено в розділі; \cite{KobayashiSchultz2008} & теорія \\
16 & \nt{SERO} задає горизонт $\tau_\gamma$ (підрозділ~\ref{sec:horizon}) & Оптогенетичне ввімкнення \nt{SERO}-нейронів дорсального шва в мишей подовжує очікування відкладеної винагороди і підвищує наполегливість у пошуку винагороди, що стала рідкісною. У людей зниження серотоніну (дієта без триптофану) робить знецінення відкладеної винагороди крутішим. Як концентрація перетворюється на $\tau_\gamma$, невідомо. & \cite{Doya2002, Miyazaki2014, Lottem2018, Schweighofer2008} & гіпотеза \\
\bottomrule
\end{longtable}}
